% Options for packages loaded elsewhere
\PassOptionsToPackage{unicode}{hyperref}
\PassOptionsToPackage{hyphens}{url}
\documentclass[
  a4paper,
]{article}
\usepackage{xcolor}
\usepackage[margin=25mm]{geometry}
\usepackage{amsmath,amssymb}
\setcounter{secnumdepth}{-\maxdimen} % remove section numbering
\usepackage{iftex}
\ifPDFTeX
  \usepackage[T1]{fontenc}
  \usepackage[utf8]{inputenc}
  \usepackage{textcomp} % provide euro and other symbols
\else % if luatex or xetex
  \usepackage{unicode-math} % this also loads fontspec
  \defaultfontfeatures{Scale=MatchLowercase}
  \defaultfontfeatures[\rmfamily]{Ligatures=TeX,Scale=1}
\fi
\usepackage{lmodern}
\ifPDFTeX\else
  % xetex/luatex font selection
\fi
% Use upquote if available, for straight quotes in verbatim environments
\IfFileExists{upquote.sty}{\usepackage{upquote}}{}
\IfFileExists{microtype.sty}{% use microtype if available
  \usepackage[]{microtype}
  \UseMicrotypeSet[protrusion]{basicmath} % disable protrusion for tt fonts
}{}
\makeatletter
\@ifundefined{KOMAClassName}{% if non-KOMA class
  \IfFileExists{parskip.sty}{%
    \usepackage{parskip}
  }{% else
    \setlength{\parindent}{0pt}
    \setlength{\parskip}{6pt plus 2pt minus 1pt}}
}{% if KOMA class
  \KOMAoptions{parskip=half}}
\makeatother
\usepackage{color}
\usepackage{fancyvrb}
\newcommand{\VerbBar}{|}
\newcommand{\VERB}{\Verb[commandchars=\\\{\}]}
\DefineVerbatimEnvironment{Highlighting}{Verbatim}{commandchars=\\\{\}}
% Add ',fontsize=\small' for more characters per line
\newenvironment{Shaded}{}{}
\newcommand{\AlertTok}[1]{\textcolor[rgb]{1.00,0.00,0.00}{\textbf{#1}}}
\newcommand{\AnnotationTok}[1]{\textcolor[rgb]{0.38,0.63,0.69}{\textbf{\textit{#1}}}}
\newcommand{\AttributeTok}[1]{\textcolor[rgb]{0.49,0.56,0.16}{#1}}
\newcommand{\BaseNTok}[1]{\textcolor[rgb]{0.25,0.63,0.44}{#1}}
\newcommand{\BuiltInTok}[1]{\textcolor[rgb]{0.00,0.50,0.00}{#1}}
\newcommand{\CharTok}[1]{\textcolor[rgb]{0.25,0.44,0.63}{#1}}
\newcommand{\CommentTok}[1]{\textcolor[rgb]{0.38,0.63,0.69}{\textit{#1}}}
\newcommand{\CommentVarTok}[1]{\textcolor[rgb]{0.38,0.63,0.69}{\textbf{\textit{#1}}}}
\newcommand{\ConstantTok}[1]{\textcolor[rgb]{0.53,0.00,0.00}{#1}}
\newcommand{\ControlFlowTok}[1]{\textcolor[rgb]{0.00,0.44,0.13}{\textbf{#1}}}
\newcommand{\DataTypeTok}[1]{\textcolor[rgb]{0.56,0.13,0.00}{#1}}
\newcommand{\DecValTok}[1]{\textcolor[rgb]{0.25,0.63,0.44}{#1}}
\newcommand{\DocumentationTok}[1]{\textcolor[rgb]{0.73,0.13,0.13}{\textit{#1}}}
\newcommand{\ErrorTok}[1]{\textcolor[rgb]{1.00,0.00,0.00}{\textbf{#1}}}
\newcommand{\ExtensionTok}[1]{#1}
\newcommand{\FloatTok}[1]{\textcolor[rgb]{0.25,0.63,0.44}{#1}}
\newcommand{\FunctionTok}[1]{\textcolor[rgb]{0.02,0.16,0.49}{#1}}
\newcommand{\ImportTok}[1]{\textcolor[rgb]{0.00,0.50,0.00}{\textbf{#1}}}
\newcommand{\InformationTok}[1]{\textcolor[rgb]{0.38,0.63,0.69}{\textbf{\textit{#1}}}}
\newcommand{\KeywordTok}[1]{\textcolor[rgb]{0.00,0.44,0.13}{\textbf{#1}}}
\newcommand{\NormalTok}[1]{#1}
\newcommand{\OperatorTok}[1]{\textcolor[rgb]{0.40,0.40,0.40}{#1}}
\newcommand{\OtherTok}[1]{\textcolor[rgb]{0.00,0.44,0.13}{#1}}
\newcommand{\PreprocessorTok}[1]{\textcolor[rgb]{0.74,0.48,0.00}{#1}}
\newcommand{\RegionMarkerTok}[1]{#1}
\newcommand{\SpecialCharTok}[1]{\textcolor[rgb]{0.25,0.44,0.63}{#1}}
\newcommand{\SpecialStringTok}[1]{\textcolor[rgb]{0.73,0.40,0.53}{#1}}
\newcommand{\StringTok}[1]{\textcolor[rgb]{0.25,0.44,0.63}{#1}}
\newcommand{\VariableTok}[1]{\textcolor[rgb]{0.10,0.09,0.49}{#1}}
\newcommand{\VerbatimStringTok}[1]{\textcolor[rgb]{0.25,0.44,0.63}{#1}}
\newcommand{\WarningTok}[1]{\textcolor[rgb]{0.38,0.63,0.69}{\textbf{\textit{#1}}}}
\usepackage{longtable,booktabs,array}
\newcounter{none} % for unnumbered tables
\usepackage{calc} % for calculating minipage widths
% Correct order of tables after \paragraph or \subparagraph
\usepackage{etoolbox}
\makeatletter
\patchcmd\longtable{\par}{\if@noskipsec\mbox{}\fi\par}{}{}
\makeatother
% Allow footnotes in longtable head/foot
\IfFileExists{footnotehyper.sty}{\usepackage{footnotehyper}}{\usepackage{footnote}}
\makesavenoteenv{longtable}
\setlength{\emergencystretch}{3em} % prevent overfull lines
\providecommand{\tightlist}{%
  \setlength{\itemsep}{0pt}\setlength{\parskip}{0pt}}
\usepackage{bookmark}
\IfFileExists{xurl.sty}{\usepackage{xurl}}{} % add URL line breaks if available
\urlstyle{same}
\hypersetup{
  pdftitle={Supplement to the Sixth Thought Experiment: Does the State Generation Change When the Normalization Reference Is Changed from Mass to a Charge Unit? - A Restricted Renormalization from G=M=c=1 to G=q\_0=c=1{,} and an Exact-Match Verification of the Initial States of the Five Stored Conditions -},
  pdfauthor={Noriaki Kihara (ORCID 0009-0004-6753-4020)},
  hidelinks,
  pdfcreator={LaTeX via pandoc}}

\title{Supplement to the Sixth Thought Experiment: Does the State
Generation Change When the Normalization Reference Is Changed from Mass
to a Charge Unit? - A Restricted Renormalization from \(G=M=c=1\) to
\(G=q_0=c=1\), and an Exact-Match Verification of the Initial States of
the Five Stored Conditions -}
\author{Noriaki Kihara (ORCID 0009-0004-6753-4020)}
\date{2026-09-27 / v1.0 / Version DOI 10.5281/zenodo.22985300}

\sloppy
\emergencystretch=3em
\begin{document}
\maketitle

\textbf{Author:} Noriaki Kihara\\
\textbf{ORCID:} 0009-0004-6753-4020\\
\textbf{Version DOI:} 10.5281/zenodo.22985300\\
\textbf{Concept DOI:} 10.5281/zenodo.22985299\\
\textbf{Version:} v1.0\\
\textbf{Date:} 2026-09-27\\
\textbf{Target paper:} The Sixth Thought Experiment v1.1, Version DOI:
10.5281/zenodo.22974632, Concept DOI: 10.5281/zenodo.22974631

\textbf{Keywords:} state closure, normalization, charge unit,
mass-to-charge ratio, initial state, charged binary, reproducibility,
deterministic state map

\begin{center}\rule{0.5\linewidth}{0.5pt}\end{center}

\section{Abstract}\label{abstract}

In the sixth thought experiment {[}6{]}, under the normalization
\(G=M=c=1\), the conservative effects of gravity and Coulomb and the
dissipative effects of GW quadrupole and EM dipole radiation of a
charged quasi-circular binary were internalized into the eight
persistent logical states

\[
Z_8=(U,P,E,H,Q,N,C,D)
\]

and a single synchronous state map \texttt{transition(z)}. There it was
confirmed, by a code audit and by numerical experiments over five
conditions, that \texttt{transition(z)} receives no external physical
argument other than the state \(z\), and that the difference of the
physical conditions is given only by the \(N,C,D\) stored in the initial
state.

In this supplement, for the same physical system, only the normalization
reference is changed from

\[
G=M=c=1
\]

to

\[
\boxed{G=q_0=c=1}.
\]

Here \(q_0\) is the reference charge unit that counts discrete charges,
and in this research

\[
q_0=\frac{e}{3}
\]

is adopted. No new dynamics, interaction term, state, update rule,
readout or approximation is added at all. The only thing changed is the
initialization that makes the initial state from the masses and the
charges.

Using equal masses

\[
m_A=m_B=\frac{M}{2}
\]

and

\[
q_A=s_A n q_0,\qquad q_B=s_B n q_0,
\qquad s_A,s_B\in\{+1,-1\},
\]

the charge-to-mass ratios of the sixth thought experiment,

\[
\lambda_A=\frac{q_A}{m_A},\qquad
\lambda_B=\frac{q_B}{m_B},
\]

become

\[
\lambda_A=s_A\frac{2nq_0}{M},\qquad
\lambda_B=s_B\frac{2nq_0}{M}.
\]

To preserve the \(n=1\) condition \(|\lambda|=0.30\) of the sixth
thought experiment, it suffices to take

\[
\boxed{\frac{M}{q_0}=\frac{20}{3}}.
\]

Using this single mass-to-charge ratio alone,
\(|\lambda|=0.30,0.90,1.20\) are regenerated for \(n=1,3,4\).

This initialization formula was implemented as an independent
initialization program and compared with the micro=0 rows of the raw
data \texttt{raw\_first\_macro\_microsteps.csv} of the five conditions
stored in the sixth thought experiment. As a result, in all five
conditions

\[
P,N,C,D,Q
\]

and the stored diagnostic columns including the initial work registers,

\texttt{micro,\ q\_index,\ P,\ N,\ C,\ D,\ k1p,\ k2p,\ k3p,\ k4p,\ dP,\ Q}

matched exactly as IEEE-754 values after reading the CSV.

In the sixth thought experiment it has already been audited that the
next state is determined from the present state alone by

\[
z_{k+1}=T(z_k).
\]

Therefore, at the point where the same initial state \(z_0\) has been
generated under the present new normalization, the whole subsequent
state sequence is also identical to the existing experiment. Re-running
the interaction runs therefore gives no new independent verification,
and is not carried out in this supplement.

The conclusion of this supplement is restricted. For the five conditions
of the sixth thought experiment, it was confirmed that even when the
normalization is changed from \(G=M=c=1\) to \(G=q_0=c=1\), the same
initial states can be generated by giving \(M/q_0=20/3\), and therefore
the existing state evolution is not changed. This is not a derivation of
a new physical law; it restates the degree of freedom that had been
absorbed into the mass reference in the sixth thought experiment,
explicitly, as a mass-to-charge ratio referred to the charge unit.

\begin{center}\rule{0.5\linewidth}{0.5pt}\end{center}

\section{1. Purpose of this
supplement}\label{purpose-of-this-supplement}

The normalization of the sixth thought experiment {[}6{]} was

\[
G=M=c=1.
\]

In this normalization the total mass \(M\) is absorbed into the unit, so
the charge condition was given to the initial state through

\[
\lambda_A=\frac{q_A}{m_A},\qquad
\lambda_B=\frac{q_B}{m_B}.
\]

This supplement asks only one thing.

\[
\boxed{
G=M=c=1
\quad\longrightarrow\quad
G=q_0=c=1
}
\]

When the reference is changed in this way and the mass \(M\) is kept
explicitly, \textbf{can exactly the same initial state as in the sixth
thought experiment be generated?}

For this question, no new orbit computation or interaction run is
needed, because it has already been audited that the generator of the
sixth thought experiment holds no physical information other than the
state.

The scope of verification of this supplement is therefore only the
following three steps.

\begin{enumerate}
\def\labelenumi{\arabic{enumi}.}
\tightlist
\item
  Confirm how the mass and charge dependence was stored into \(C,D\) in
  the initialization code of the sixth thought experiment.
\item
  Replace only that initialization part by a representation with
  \(G=q_0=c=1\) and variable \(M\).
\item
  Compare the initial rows generated by the new initializer directly
  with the stored raw initial rows of the sixth thought experiment.
\end{enumerate}

The update map \texttt{transition(z)}, the 11-phase machine, the RK4
work registers, the readout, the numerical integration conditions and
the existing raw orbit data are not changed.

\begin{center}\rule{0.5\linewidth}{0.5pt}\end{center}

\section{2. What is already established in the sixth thought
experiment}\label{what-is-already-established-in-the-sixth-thought-experiment}

\subsection{2.1 State closure}\label{state-closure}

In the sixth thought experiment the persistent logical states were taken
to be

\[
Z_8=(U,P,E,H,Q,N,C,D),
\]

and a full microstate of 41 real components was constructed, including
the RK4 work registers and the 11-phase one-hot phase {[}6{]}.

What matters is that the sole generating map is

\begin{Shaded}
\begin{Highlighting}[]
\NormalTok{transition(z)}
\end{Highlighting}
\end{Shaded}

and that it receives nothing other than the present state \(z\) as an
argument.

The generating rule can therefore be written abstractly as

\[
\boxed{z_{k+1}=T(z_k)}.
\]

\subsection{2.2 The physical condition is stored in the initial
state}\label{the-physical-condition-is-stored-in-the-initial-state}

In the sixth thought experiment the charge dependence was organized as

\[
C=1-\lambda_A\lambda_B,
\]

\[
D=(\lambda_A-\lambda_B)^2,
\]

where

\[
\lambda_A=\frac{q_A}{m_A},\qquad
\lambda_B=\frac{q_B}{m_B}.
\]

In the implementation, the difference between the five conditions is
only the \(C,D\) stored in the initial state, and the same
\texttt{transition(z)} was used for all conditions {[}6{]}.

In this supplement, therefore, it suffices to regenerate \(C,D\) as the
same values.

\begin{center}\rule{0.5\linewidth}{0.5pt}\end{center}

\section{3. Change of the
normalization}\label{change-of-the-normalization}

\subsection{3.1 The old normalization}\label{the-old-normalization}

The sixth thought experiment adopted

\[
G=M=c=1.
\]

Here \(M\) is the total mass of the binary. Under the equal-mass
condition,

\[
m_A=m_B=\frac{M}{2}=\frac12.
\]

The charge-to-mass ratios used in the five conditions are

\[
(\lambda_A,\lambda_B)
=
(+0.30,-0.30),
(+0.30,+0.30),
(+0.90,-0.90),
(+0.90,+0.90),
(+1.20,-1.20).
\]

\subsection{3.2 The new normalization}\label{the-new-normalization}

In this supplement,

\[
\boxed{G=q_0=c=1}.
\]

\(q_0\) is the reference charge unit,

\[
q_0=\frac{e}{3}.
\]

The charges are expressed as integer multiples,

\[
q_A=s_A n q_0,
\qquad
q_B=s_B n q_0,
\]

where

\[
s_A,s_B\in\{+1,-1\},\qquad n\in\mathbb N.
\]

On the other hand, the mass is not absorbed into the unit:

\[
m_A=m_B=\frac{M}{2}.
\]

Then

\[
\lambda_A
=
\frac{s_A n q_0}{M/2}
=
s_A\frac{2nq_0}{M},
\]

\[
\lambda_B
=
s_B\frac{2nq_0}{M}.
\]

The continuous degree of freedom remaining after the change of
normalization is therefore

\[
\boxed{\frac{M}{q_0}}.
\]

\begin{center}\rule{0.5\linewidth}{0.5pt}\end{center}

\section{4. The mass-to-charge ratio that regenerates the old initial
conditions}\label{the-mass-to-charge-ratio-that-regenerates-the-old-initial-conditions}

In the \(n=1\) condition of the sixth thought experiment,

\[
|\lambda|=0.30=\frac{3}{10}.
\]

In the new normalization,

\[
|\lambda|
=
\frac{2q_0}{M}.
\]

Hence

\[
\frac{2q_0}{M}=\frac{3}{10},
\]

that is,

\[
\boxed{\frac{M}{q_0}=\frac{20}{3}}.
\]

In units with \(q_0=1\),

\[
\boxed{M=\frac{20}{3}},
\]

\[
\boxed{m_A=m_B=\frac{10}{3}}.
\]

Fixing this single value, for a general integer \(n\)

\[
|\lambda_n|
=
\frac{2n}{20/3}
=
\frac{3n}{10}.
\]

Therefore

\[
n=1\Rightarrow |\lambda|=0.30,
\]

\[
n=3\Rightarrow |\lambda|=0.90,
\]

\[
n=4\Rightarrow |\lambda|=1.20.
\]

The five conditions of the sixth thought experiment are regenerated as
they are from a single \(M/q_0\).

\begin{center}\rule{0.5\linewidth}{0.5pt}\end{center}

\section{\texorpdfstring{5. Regeneration of the initial states
\(C,D\)}{5. Regeneration of the initial states C,D}}\label{regeneration-of-the-initial-states-cd}

Expressing equal and opposite signs by \(s_A,s_B\),

\[
\lambda_A=s_A\frac{3n}{10},
\qquad
\lambda_B=s_B\frac{3n}{10}.
\]

Therefore

\[
C
=
1-\lambda_A\lambda_B
=
1-s_As_B\frac{9n^2}{100},
\]

\[
D
=
(\lambda_A-\lambda_B)^2
=
\frac{9n^2}{100}(s_A-s_B)^2.
\]

For opposite signs \(s_A=-s_B\),

\[
C=1+\frac{9n^2}{100},
\]

\[
D=\frac{36n^2}{100}.
\]

For equal signs \(s_A=s_B\),

\[
C=1-\frac{9n^2}{100},
\]

\[
D=0.
\]

Hence the five conditions are

{\def\LTcaptype{none} % do not increment counter
\begin{longtable}[]{@{}lrllrr@{}}
\toprule\noalign{}
case & \(n\) & \((s_A,s_B)\) & \((\lambda_A,\lambda_B)\) & \(C\) &
\(D\) \\
\midrule\noalign{}
\endhead
\bottomrule\noalign{}
\endlastfoot
n1\_attractive & 1 & \((+,-)\) & \((+0.30,-0.30)\) & 1.09 & 0.36 \\
n1\_repulsive & 1 & \((+,+)\) & \((+0.30,+0.30)\) & 0.91 & 0.00 \\
n3\_attractive & 3 & \((+,-)\) & \((+0.90,-0.90)\) & 1.81 & 3.24 \\
n3\_repulsive & 3 & \((+,+)\) & \((+0.90,+0.90)\) & 0.19 & 0.00 \\
n4\_attractive & 4 & \((+,-)\) & \((+1.20,-1.20)\) & 2.44 & 5.76 \\
\end{longtable}
}

which coincide with the initial values of the sixth thought experiment.

\begin{center}\rule{0.5\linewidth}{0.5pt}\end{center}

\section{6. Code verification}\label{code-verification}

\subsection{6.1 Scope of the change}\label{scope-of-the-change}

In the existing code of the sixth thought experiment, the case table
gave \(\lambda_A,\lambda_B,C,D\) as in

\begin{Shaded}
\begin{Highlighting}[]
\NormalTok{(}\StringTok{"n1\_attractive"}\NormalTok{, }\FloatTok{0.30}\NormalTok{, }\OperatorTok{{-}}\FloatTok{0.30}\NormalTok{, }\FloatTok{1.09}\NormalTok{, }\FloatTok{0.36}\NormalTok{)}
\end{Highlighting}
\end{Shaded}

and the final \(C,D\) were passed to \texttt{init\_state()}.

In this supplement \texttt{transition(z)} is not changed. A new,
independent initialization module

\begin{Shaded}
\begin{Highlighting}[]
\NormalTok{paper6\_init\_G\_q0\_c1.py}
\end{Highlighting}
\end{Shaded}

was created, which takes as input

\[
(q_0,M,n,s_A,s_B)
\]

generates

\[
\lambda_A,\lambda_B,C,D,
\]

and from those values makes the same initial diagnostic row as the sixth
thought experiment.

\subsection{6.2 Test stub}\label{test-stub}

A test stub that generates the five conditions at once,

\begin{Shaded}
\begin{Highlighting}[]
\NormalTok{test\_stub\_generate\_initial\_rows.py}
\end{Highlighting}
\end{Shaded}

was created. Its output is

\begin{Shaded}
\begin{Highlighting}[]
\NormalTok{generated\_initial\_rows\_G\_q0\_c1.csv}
\end{Highlighting}
\end{Shaded}

As the object of comparison, the micro=0 row was extracted from each

\begin{Shaded}
\begin{Highlighting}[]
\NormalTok{raw\_first\_macro\_microsteps.csv}
\end{Highlighting}
\end{Shaded}

of the five stored cases of the sixth thought experiment. The result of
the extraction was fixed as

\begin{Shaded}
\begin{Highlighting}[]
\NormalTok{original\_initial\_rows\_paper6.csv}
\end{Highlighting}
\end{Shaded}

\subsection{6.3 Method of comparison}\label{method-of-comparison}

With the comparison program

\begin{Shaded}
\begin{Highlighting}[]
\NormalTok{compare\_initial\_rows.py}
\end{Highlighting}
\end{Shaded}

the following 12 columns were compared.

\begin{Shaded}
\begin{Highlighting}[]
\NormalTok{micro, q\_index, P, N, C, D,}
\NormalTok{k1p, k2p, k3p, k4p, dP, Q}
\end{Highlighting}
\end{Shaded}

The comparison was made as an exact match of IEEE-754 floats after
reading the CSV as numbers.

Note that \(M/q_0=20/3\) and \(\lambda=3n/10\) were constructed as
rational numbers and converted to float at the end. This is an
implementation measure to reproduce decimal notations such as
\texttt{1.81} as the same IEEE-754 values as the old raw data; it is not
a change of the physical model.

\begin{center}\rule{0.5\linewidth}{0.5pt}\end{center}

\section{7. Verification results}\label{verification-results}

The result was PASS in all five conditions.

{\def\LTcaptype{none} % do not increment counter
\begin{longtable}[]{@{}lrrrrrl@{}}
\toprule\noalign{}
case & \(P\) & \(N\) & \(C\) & \(D\) & \(Q\) & result \\
\midrule\noalign{}
\endhead
\bottomrule\noalign{}
\endlastfoot
n1\_attractive & 50.0 & 0.25 & 1.09 & 0.36 & 1.0 & PASS \\
n1\_repulsive & 50.0 & 0.25 & 0.91 & 0.00 & 1.0 & PASS \\
n3\_attractive & 50.0 & 0.25 & 1.81 & 3.24 & 1.0 & PASS \\
n3\_repulsive & 50.0 & 0.25 & 0.19 & 0.00 & 1.0 & PASS \\
n4\_attractive & 50.0 & 0.25 & 2.44 & 5.76 & 1.0 & PASS \\
\end{longtable}
}

For all cases,

\[
\boxed{
\Delta P=
\Delta N=
\Delta C=
\Delta D=
\Delta Q=0
}
\]

and for the initial work registers as well,

\[
\boxed{
\Delta k_{1p}=
\Delta k_{2p}=
\Delta k_{3p}=
\Delta k_{4p}=
\Delta dP=0
}
\]

Furthermore,

\[
\Delta \mathrm{micro}=0,
\qquad
\Delta q_{\mathrm{index}}=0.
\]

The initial rows observed in the stored raw data therefore matched
exactly in all 12 columns.

The machine-readable detailed results were saved in

\begin{Shaded}
\begin{Highlighting}[]
\NormalTok{initial\_row\_verification\_result.json}
\end{Highlighting}
\end{Shaded}

\begin{center}\rule{0.5\linewidth}{0.5pt}\end{center}

\section{8. Why the interaction runs are not
re-executed}\label{why-the-interaction-runs-are-not-re-executed}

This is an important point of this supplement.

In the sixth thought experiment it has already been audited that the
generating map generates the next state from the state \(z_k\) alone, as

\[
z_{k+1}=T(z_k),
\]

and does not refer to \(G,M,q_A,q_B\), the case name or the like as
external arguments {[}6{]}.

Let the initial state under the old normalization be
\(z_0^{\mathrm{old}}\) and the initial state under the new normalization
be \(z_0^{\mathrm{new}}\).

By the present verification,

\[
\boxed{
z_0^{\mathrm{new}}=z_0^{\mathrm{old}}
}
\]

was confirmed for the initialization part that is affected by the change
of normalization.

Since the same map \(T\) is used,

\[
z_1^{\mathrm{new}}
=T(z_0^{\mathrm{new}})
=T(z_0^{\mathrm{old}})
=z_1^{\mathrm{old}}.
\]

Likewise, by induction,

\[
\boxed{
z_k^{\mathrm{new}}=z_k^{\mathrm{old}}
\qquad (k=0,1,2,\ldots)
}
\]

Therefore, re-executing the interaction runs for the five existing
conditions would only regenerate the same deterministic state sequence.
That is not a new independent verification.

For this reason, in this supplement \textbf{only the verification of the
initialization equivalence is carried out, and the orbit runs are not
re-executed}.

\begin{center}\rule{0.5\linewidth}{0.5pt}\end{center}

\section{9. Interpretation}\label{interpretation}

\subsection{\texorpdfstring{9.1 \(M=1\) gave no additional information
to the dynamics of the present five
conditions}{9.1 M=1 gave no additional information to the dynamics of the present five conditions}}\label{m1-gave-no-additional-information-to-the-dynamics-of-the-present-five-conditions}

The sixth thought experiment used \(M=1\) for the normalization. In this
supplement the charge unit \(q_0\) was set to 1 and, instead,

\[
\frac{M}{q_0}
\]

was kept.

As a result, with

\[
\frac{M}{q_0}=\frac{20}{3}
\]

the same initial states as in the sixth thought experiment were
obtained.

Therefore, at least for the five conditions verified here, the choice
\(M=1\) itself did not give additional independent information to the
state evolution. The necessary information can be re-expressed as a
dimensionless mass-to-charge ratio at the time of initialization.

\subsection{9.2 Discrete charge and continuous mass can be described
separately}\label{discrete-charge-and-continuous-mass-can-be-described-separately}

In the new representation, the charge side is expressed by the integer
\(n\) through

\[
q=nq_0,
\]

and the mass side can be kept as

\[
M/q_0.
\]

In the present five conditions, fixing the single

\[
M/q_0=20/3
\]

and only changing \(n=1,3,4\) regenerated the old conditions.

However, this supplement does not derive a new quantization rule from
this discreteness. What was confirmed here is only the fact that
\textbf{the existing conditions of the sixth thought experiment can be
rewritten, without loss, into an initialization referred to the charge
unit}.

\begin{center}\rule{0.5\linewidth}{0.5pt}\end{center}

\section{10. What this supplement does not
show}\label{what-this-supplement-does-not-show}

This supplement does not claim the following.

\begin{itemize}
\tightlist
\item
  That \(q_0=e/3\) has been derived from this theory.
\item
  That \(M/q_0=20/3\) is a universal physical constant.
\item
  That the origin of mass quantization or charge quantization has been
  explained.
\item
  Validity for general charged binaries beyond the range of
  approximation of the sixth thought experiment.
\item
  That a fundamental unification of gravity and electromagnetism has
  been derived.
\item
  That a new orbit or a new radiation phenomenon has been discovered.
\end{itemize}

The object of verification of this supplement is only the
\textbf{equivalence of the initialization normalization} of the five
conditions already executed in the sixth thought experiment.

\begin{center}\rule{0.5\linewidth}{0.5pt}\end{center}

\section{11. Files for reproduction}\label{files-for-reproduction}

The files needed to reproduce this supplement were saved in the folder
(the folder names in the repository are in Japanese; transliterated
here)

\begin{Shaded}
\begin{Highlighting}[]
\NormalTok{04\_supplement\_Gq0c1\_initialization\_equivalence\_verification\_20260927/}
\end{Highlighting}
\end{Shaded}

inside the storage folder of the sixth thought experiment.

The main files are as follows.

\begin{Shaded}
\begin{Highlighting}[]
\NormalTok{paper6\_init\_G\_q0\_c1.py}
\end{Highlighting}
\end{Shaded}

The restricted initialization implementation that generates the initial
states of the five conditions from \(G=q_0=c=1\), \(M/q_0=20/3\).

\begin{Shaded}
\begin{Highlighting}[]
\NormalTok{test\_stub\_generate\_initial\_rows.py}
\end{Highlighting}
\end{Shaded}

The test stub that generates the initial rows of the five conditions.

\begin{Shaded}
\begin{Highlighting}[]
\NormalTok{generated\_initial\_rows\_G\_q0\_c1.csv}
\end{Highlighting}
\end{Shaded}

The output of the new initializer.

\begin{Shaded}
\begin{Highlighting}[]
\NormalTok{original\_initial\_rows\_paper6.csv}
\end{Highlighting}
\end{Shaded}

The object of comparison extracted from the stored raw data of the sixth
thought experiment.

\begin{Shaded}
\begin{Highlighting}[]
\NormalTok{compare\_initial\_rows.py}
\end{Highlighting}
\end{Shaded}

The program verifying the exact match of the new and old initial rows.

\begin{Shaded}
\begin{Highlighting}[]
\NormalTok{initial\_row\_verification\_result.json}
\end{Highlighting}
\end{Shaded}

The comparison results of all fields.

\begin{Shaded}
\begin{Highlighting}[]
\NormalTok{initial\_row\_verification\_result\_ja.md}
\end{Highlighting}
\end{Shaded}

A short summary of the verification results in Japanese.

\begin{center}\rule{0.5\linewidth}{0.5pt}\end{center}

\section{12. Conclusion}\label{conclusion}

The normalization of the sixth thought experiment,

\[
G=M=c=1,
\]

was changed to

\[
G=q_0=c=1,
\]

and the mass was kept as an explicit degree of freedom.

For the equal-mass binary and the charge unit \(q_0\),

\[
\lambda_{A,B}=s_{A,B}\frac{2nq_0}{M},
\]

and from the reference value \(|\lambda|=0.30\) of the sixth thought
experiment,

\[
\boxed{\frac{M}{q_0}=\frac{20}{3}}
\]

was obtained.

Initializing the five conditions of \(n=1,3,4\) with this single ratio,
they matched exactly, in all compared columns, the stored initial raw
rows of the sixth thought experiment.

Therefore, within the present scope of verification,

\[
\boxed{
G=M=c=1
\quad\text{and}\quad
G=q_0=c=1,\;M/q_0=20/3
}
\]

are \textbf{two normalization representations that generate the same
internal initial state}.

Under the state closure

\[
z_{k+1}=T(z_k)
\]

already established in the sixth thought experiment, if the initial
state is the same then all subsequent states are also the same. For this
reason, re-execution of the interaction runs is unnecessary.

By this supplement it was confirmed that moving the degree of freedom
that had been absorbed into the mass unit in the sixth thought
experiment into the single mass-to-charge ratio

\[
\boxed{M/q_0}
\]

referred to the charge unit does not affect the generation of the
physical states of the five existing conditions.

\begin{center}\rule{0.5\linewidth}{0.5pt}\end{center}

\section{References}\label{references}

\subsection{This series}\label{this-series}

{[}0{]} Noriaki Kihara, \textbf{Designing a Research Project to Explore
the Emergence of Spacetime, Particles, and Interactions from States and
Relations: A Framework of Preliminary Design and Feasibility
Verification for Foundational-Physics Model Exploration under
Uncertainty}, v1.0 (2026-09-20). Concept DOI: 10.5281/zenodo.22851944;
Version DOI: 10.5281/zenodo.22851945.

{[}1{]} Noriaki Kihara, \textbf{The First Thought Experiment: Tracing
the Equivalence Principle Down to Two Nameless Degrees of Freedom -
Deriving the Sign of the Sum of Squares, the Curvature Radius, and Three
Systems (Inertial, Uniformly Accelerated, Rotating) from the
Conservation of an Area Readout}, v1.0 (2026-09-20). Concept DOI:
10.5281/zenodo.22857952; Version DOI: 10.5281/zenodo.22857953.

{[}2{]} Noriaki Kihara, \textbf{The Second Thought Experiment: Deriving
the Coulomb-Type Inverse-Square Force from a Product Readout and an Area
Clock - Keeping Two Values and a Linear Law, Obtaining the Focus,
Attraction and Repulsion and Neutrality, and Fixing What Two Values
Cannot Write}, v1.2 (2026-09-21). Concept DOI: 10.5281/zenodo.22867336;
Version DOI: 10.5281/zenodo.22876596.

{[}3{]} Noriaki Kihara, \textbf{The Third Thought Experiment: Reading
Out the State of Two Bodies a, b from the Orbit of Two Celestial Bodies
a, b - Obtaining the Mass Ratio, the Relative Distance and the Relative
Time from Orbit and Gravitational-Wave Readouts, with Two Values of
Undetermined Meaning and the Knowledge \(G=c=M=1\) Alone}, v1.5
(2026-09-23). Concept DOI: 10.5281/zenodo.22909660; Version DOI:
10.5281/zenodo.22909661.

{[}4{]} Noriaki Kihara, \textbf{The Fourth Thought Experiment: From a
Fixed Action \(S\) to a State-Dependent Action \(S_n\) - Can the
Relativistic Gravitational Orbits Computed Externally in the Third
Thought Experiment Be Regenerated from the Sequential Interaction
\(X_{n+1}=S_nX_n\)?}, v1.0 (2026-09-23). Concept DOI:
10.5281/zenodo.22919787; Version DOI: 10.5281/zenodo.22919788.

{[}5{]} Noriaki Kihara, \textbf{The Fifth Thought Experiment: Was There
No Hidden State? - A Self-Audit of the Two-State Orbit Generation of the
Fourth Thought Experiment, Tracking the State Closure from Five States
to the Sixth Persistent State \(\mathcal N=\nu\) -}, v2.1 (2026-09-26).
Concept DOI: 10.5281/zenodo.22973086; Version DOI:
10.5281/zenodo.22973087.

{[}6{]} Noriaki Kihara, \textbf{The Sixth Thought Experiment: Can Two
Kinds of Long-Range Interaction Be Internalized into a Single Closed
State Map? - Reconstructing Gravity, Coulomb, GW-Quadrupole and
EM-Dipole Radiation of a Charged Quasi-Circular Binary with Eight
Persistent States, and Verifying the Transferability of the Same
Transition over Five Conditions -}, v1.1 (2026-09-26). Concept DOI:
10.5281/zenodo.22974631; Version DOI: 10.5281/zenodo.22974632.

\subsection{External}\label{external}

{[}7{]} P. C. Peters and J. Mathews, ``Gravitational Radiation from
Point Masses in a Keplerian Orbit,'' \emph{Physical Review}
\textbf{131}, 435-440 (1963). DOI: 10.1103/PhysRev.131.435.

{[}8{]} P. C. Peters, ``Gravitational Radiation and the Motion of Two
Point Masses,'' \emph{Physical Review} \textbf{136}, B1224-B1232 (1964).
DOI: 10.1103/PhysRev.136.B1224.

{[}9{]} L. Liu, O. Christiansen, Z.-K. Guo, R.-G. Cai, and S. P. Kim,
``Gravitational and electromagnetic radiation from binary black holes
with electric and magnetic charges: Circular orbits on a cone,''
\emph{Physical Review D} \textbf{102}, 103520 (2020). DOI:
10.1103/PhysRevD.102.103520.

{[}10{]} C. A. Benavides-Gallego and W.-B. Han, ``Gravitational Waves
and Electromagnetic Radiation from Charged Black Hole Binaries,''
\emph{Symmetry} \textbf{15}, 537 (2023). DOI: 10.3390/sym15020537.

{[}11{]} Z.-H. Zhang, T. Liu, S. Zhang, and Z.-K. Guo, ``Post-Newtonian
dynamics of charged compact binaries,'' \emph{Physical Review D}
\textbf{114}, 044088 (2026). DOI: 10.1103/4x8q-2kzm.

{[}12{]} S. Verma et al., ``Post-newtonian dynamics of radiating
charges: canonical formulation and binary inspiral laws,''
\emph{European Physical Journal C} \textbf{86}, 830 (2026). DOI:
10.1140/epjc/s10052-026-16076-2.

\begin{center}\rule{0.5\linewidth}{0.5pt}\end{center}

\subsection{Reproducibility note}\label{reproducibility-note}

This supplement does not change the dynamics of the existing sixth
thought experiment. The new code used for the verification, the
generated initial rows, the initial rows extracted from the old raw
data, the comparison code and the comparison results were placed in the
same storage folder described in Section 11 of the main text. For the
interaction-run data, the existing raw data of the sixth thought
experiment are taken as the object of verification as they are, and are
not regenerated.

\end{document}
